\documentclass[11pt]{article}
\usepackage[a4paper,margin=1in]{geometry}
\usepackage{amsmath,amssymb,amsthm,mathtools}

\newtheorem{theorem}{Theorem}[section]
\newtheorem{lemma}[theorem]{Lemma}
\newtheorem{proposition}[theorem]{Proposition}
\newtheorem{conjecture}[theorem]{Conjecture}
\newtheorem{remark}[theorem]{Remark}

\newcommand{\Om}{\Omega}
\newcommand{\R}{\mathbb R}
\newcommand{\Sph}{\mathbb S}
\newcommand{\eps}{\varepsilon}
\newcommand{\cA}{\mathcal A}
\newcommand{\dH}{\,d\mathcal H}

\title{Plan 3 Single-Level Recheck and the Planar Annular Bridge}
\author{}
\date{May 18, 2026}

\begin{document}
\maketitle

\section{Outcome}

The sound Plan 3 information is more useful than the old graph route, but it
does not yet close a new proof.

\begin{itemize}
\item The Chebyshev good-level extraction, the deficit transfer
from \(\Om\) to \(E_{\hat t}\), and the final superlevel-to-domain
Fraenkel transfer are sound quantitative steps.
\item The recently proved small-amplitude radial-graph theorem is also sound:
for star-shaped radial graphs in a thin annulus it gives
\[
   E(D)-E(B)\ge c_n\|R-1\|_{L^2(\Sph^{n-1})}^2\ge c_n'\cA(D)^2
\]
with no slope, curvature, \(C^1\), or perimeter assumption.
\item The single-level Hessian/Weinberger ``Chain B'' currently recorded in
\texttt{proof-step2.md} and \texttt{verify-chainB.md} is not sound.  The
exact false point is the unweighted Hessian domination
\[
   \delta_T(D)\ge c_n |D|^{-1}\int_D \left|D^2u_D+\frac1n I\right|^2 .
\]
This has the wrong derivative count and fails on smooth nearly spherical
high-frequency perturbations.
\item In dimension \(2\), the most promising non-Plan-1 route is therefore:
good level \(\Rightarrow\) centered small annular containment, followed by a
new annular-set stability theorem modulo translations.  Neither implication
is currently proved in the repo.  The first is only recorded as a
reconstructed status note; the second is not present at all.
\end{itemize}

\section{The precise overclaim: unweighted Hessian control}

The claimed Chain B uses the following statement, labelled \((T)\) in
\texttt{Plan 3/C-assembled-proof-steps/proof-step2.md}:
\[
   \delta_T(D)\gtrsim_n |D|^{-1}\int_D
       \left|D^2u_D+\frac1n I\right|^2 .
   \tag{T}
\]
The verification note \texttt{Plan 3/B-routeB-weinberger/verify-chainB.md}
then treats this as a regularity-free Weinberger quantitative inequality.
This is the point where compactness/regularity was not eliminated; a much
stronger and false Hessian statement was substituted for the actual
Saint--Venant deficit.

\begin{proposition}[High-frequency obstruction to \((T)\)]
There is no dimensional \(c_n>0\) such that \((T)\) holds for all smooth
volume-normalized near-balls.
\end{proposition}

\begin{proof}
Let \(Y_k\) be an \(L^2(\Sph^{n-1})\)-normalized spherical harmonic of degree
\(k\ge2\).  Consider the smooth radial graph
\[
   D_{\eps,k}=\{r< R_{\eps,k}(\theta)\},\qquad
   R_{\eps,k}=1+\eps Y_k+O_k(\eps^2),
\]
where the \(O_k(\eps^2)\) constant mode enforces volume and no first-mode
correction is needed because \(k\ge2\).  The expansion below is obtained by
the elementary ball linearization, equivalently by the exact radial-graph
identity in \texttt{H12\_ASYMMETRY\_DIRECT\_PROOF.tex}; no BDV theorem is
being used.

Set
\[
   H_k(r,\theta)=r^kY_k(\theta).
\]
The first variation of \(u_D-u_B\) is \(\eps H_k/n\).  The torsion deficit
has the Steklov, half-derivative scale
\[
   \delta_T(D_{\eps,k})
      = \frac{\eps^2}{2n^2}(k-1)+O_k(\eps^3).
   \tag{2.1}
\]
Indeed the harmonic energy contributes
\(\frac12(\eps^2/n^2)\int_{B_1}|\nabla H_k|^2
=\frac{\eps^2}{2n^2}k\), while the volume correction in the moment term
contributes \(-\eps^2/(2n^2)\).

By contrast, the unweighted Hessian defect has two more powers of frequency.
Using
\[
   \Delta |\nabla H_k|^2=2|D^2H_k|^2,
\]
and
\[
   |\nabla H_k|^2\big|_{\partial B_1}
      = k^2Y_k^2+|\nabla_\tau Y_k|^2,\qquad
   \int_{\Sph^{n-1}}|\nabla_\tau Y_k|^2
      = k(k+n-2),
\]
we get
\[
   \int_{B_1}|D^2H_k|^2
       = (k-1)\bigl(k^2+k(k+n-2)\bigr)
       = (k-1)k(2k+n-2).
   \tag{2.2}
\]
Therefore
\[
   \int_{D_{\eps,k}}\left|D^2u_{D_{\eps,k}}+\frac1nI\right|^2
      = \frac{\eps^2}{n^2}(k-1)k(2k+n-2)+O_k(\eps^3).
   \tag{2.3}
\]
Taking \(\eps\downarrow0\) for fixed \(k\), then \(k\to\infty\), gives
\[
   \frac{\delta_T(D_{\eps,k})}
        {\int_{D_{\eps,k}}|D^2u_{D_{\eps,k}}+\frac1nI|^2}
      \longrightarrow
      \frac{1}{2k(2k+n-2)}\longrightarrow0.
\]
Thus no dimensional lower bound \((T)\) can be true.
\end{proof}

The weighted Hessian variant is also false at the sharp scale, but less
violently: \(\int_{B_1}u_B|D^2H_k|^2=k(k-1)/n\), so it sees an \(H^1\)
boundary scale.  The unweighted Chain B sees an even stronger scale and
therefore cannot be the missing Plan 3 closure.

\section{What remains sound}

The following pieces survive this recheck.

\begin{enumerate}
\item \textbf{Step 1.}  The good-level extraction gives a Sard-regular
level \(E_{\hat t}\) with small \(D_H\)-integrand and
\[
   |\Om\setminus E_{\hat t}|\le C_{n,\theta}\,\delta_T(\Om).
\]
The important patch is that the argument is made in the volume variable
\(s=m(t)\), where the layer size is exactly linear.

\item \textbf{Step 3.}  The function \(u_\Om-\hat t\) is exactly the torsion
function of \(E_{\hat t}\).  The profile computation gives
\[
   \delta_T(E_{\hat t})\le
   \frac{1}{1-|\Om\setminus E_{\hat t}|/|\Om|}\,\delta_T(\Om),
\]
hence \(\delta_T(E_{\hat t})\le 2\delta_T(\Om)\) in the small-deficit
regime.

\item \textbf{Step 5.}  For \(E_{\hat t}\subset \Om\),
\[
   \cA(\Om)\le \cA(E_{\hat t})+
       2\frac{|\Om\setminus E_{\hat t}|}{|\Om|}.
\]
After squaring, the transfer error is \(O(\delta_T(\Om)^2)\).

\item \textbf{Small-amplitude radial graphs.}  The direct radial-graph proof
closes the desired \(A^2\) theorem in a thin annulus for genuine radial
graphs.  This is stronger than the old graph route in one direction: it
does not need any slope or curvature bound.  It still needs one radial
crossing per direction.
\end{enumerate}

\section{The planar route that is still plausible}

The \(n=2\) note \texttt{n2-graph-route-OBSTRUCTED.md} records the right
positive target: near-Serrin data on a planar level should perhaps give only
\(L^\infty\)/inner--outer radius control, not \(C^1\), \(H^1\), or curvature
control.  But this is a reconstructed verdict, not a written proof in the
repo.  With the new radial-graph theorem, such a weak conclusion would no
longer be obviously too weak.  The graph theorem only needs small amplitude,
not small slope.

There are therefore two separate missing statements.  First, one needs the
annular entry lemma.

\begin{conjecture}[Planar annular entry from the good level]
Let \(E_{\hat t}\subset\R^2\) be the Sard-regular level obtained by the
Plan 3 Step 1 extraction, with normalized area and \(D_H(\hat t)\le\theta\).
There exist a center \(x_0\) and a modulus \(\eta(\theta)\to0\), with
computable constants, such that
\[
   B_{1-\eta(\theta)}(x_0)\subset E_{\hat t}
      \subset B_{1+\eta(\theta)}(x_0).
   \tag{AE}
\]
\end{conjecture}

This would be weaker than graph entry and is exactly why dimension \(2\) is
still worth pursuing.  But (AE) is not proved in the current notes.

There is one neutral-mode warning before formulating the second statement:
the fixed-center inequality
\[
   E(D)-E(B_1)\ge c\,|D\Delta B_1|^2
\]
is false under only \(B_{1-\eta}\subset D\subset B_{1+\eta}\).  A translated
disk \(D=B_1(a)\), \(0<|a|\le\eta\), satisfies the annular containment and
has \(E(D)=E(B_1)\), but \(|D\Delta B_1|>0\).  The bridge must either impose
a centering condition eliminating the first Fourier mode, or state the
right side with the Fraenkel infimum over translated balls.

Second, even if (AE) were available with the correct center, annular
containment is not the same as being a radial graph.  A level set may have
several crossings on one ray, holes/fjords inside the annulus, or same-ray
cancellation between an inward missing piece and an outward excess piece.
Thus the second missing statement is the following.

\begin{conjecture}[Planar annular bridge, centered form]
There exist computable constants \(\eta_0,c_0>0\) such that, if
\[
   |D|=\pi,\qquad B_{1-\eta}\subset D\subset B_{1+\eta},
   \qquad \int_D x\,dx=0,\qquad 0<\eta\le\eta_0,
\]
then
\[
   E(D)-E(B_1)\ge c_0\,|D\Delta B_1|^2 .
   \tag{AB}
\]
Equivalently, in translation-invariant language, one may replace the
right-hand side by \(c_0\inf_a |D\Delta B_1(a)|^2\), provided the annular
entry lemma also supplies a compatible center.
\end{conjecture}

If both (AE) and (AB) are proved with compatible centers, the planar
single-level route closes without BDV: Step 1 picks \(E_{\hat t}\) with a
small \(D_H\)-level; the planar annular-entry lemma gives
\[
   B_{1-\eta(\theta)}\subset E_{\hat t}^{\sharp}\subset B_{1+\eta(\theta)},
   \qquad \eta(\theta)\lesssim \theta^{1/2};
\]
choose the fixed \(\theta\) so that \(\eta(\theta)\le\eta_0\); Step 3 gives
\(\delta_T(E_{\hat t})\le C\delta_T(\Om)\); (AB) gives
\(\cA(E_{\hat t})^2\le C\delta_T(E_{\hat t})\); Step 5 transfers this back
to \(\Om\).

\section{Attempted proof of the annular bridge}

There is a clean identity for arbitrary planar annular sets which identifies
the real issue.  Let \(D\subset\R^2\) be smooth, \(|D|=\pi\), and
\[
   u_0(x)=\frac{1-|x|^2}{4},\qquad z=u_D-u_0.
\]
Then \(z\) is harmonic in \(D\) and \(z=-u_0\) on \(\partial D\).  Expanding
the torsion energy around \(u_0\) gives
\[
   E(D)-E(B_1)
     = \frac12\int_D|\nabla z|^2\,dx
       -\frac18\left(\int_D |x|^2\,dx-\int_{B_1}|x|^2\,dx\right).
   \tag{4.1}
\]
For \(B_{1-\eta}\subset D\subset B_{1+\eta}\), the moment term is
\[
   M_2(D):=
   \int_{D\setminus B_1}(|x|^2-1)\,dx
   +\int_{B_1\setminus D}(1-|x|^2)\,dx \ge0,
\]
so (4.1) reads
\[
   E(D)-E(B_1)=\frac12\int_D|\nabla z|^2-\frac18 M_2(D).
   \tag{4.2}
\]

The geometric side is also clean once the translation mode has been fixed.
With
\[
   M_1(D):=\int_{D\Delta B_1}\bigl||x|-1\bigr|\,dx,
\]
the annular bounds imply \(M_2(D)\simeq M_1(D)\).  The slice
Cauchy--Schwarz inequality already present in Plan 2 gives, in this centered
annular cap,
\[
   |D\Delta B_1|^2 \le C\,M_1(D)
   \tag{4.3}
\]
with a computable absolute constant when \(\eta\le\eta_0\).
Thus a very natural sufficient estimate for (AB) is
\[
   E(D)-E(B_1)\ge c\,M_1(D).
   \tag{4.4}
\]

For radial graphs, (4.4) is exactly what the direct graph proof establishes
after removing the neutral volume and translation modes.  For general
annular sets, (4.4) is not yet proved in the repo.  The missing analytic
lemma is a capacitary radial-moment coercivity estimate:
\[
   \int_D|\nabla z|^2
      \ge \left(\frac14+c\right)M_2(D)
        \quad\hbox{modulo the volume and translation neutral modes.}
   \tag{4.5}
\]

The hard part is not high frequency by itself, nor translation.  Translation
is the first harmonic and must be quotiented out or fixed by centering.
High-frequency radial corrugations are already handled by the slope-free
radial-graph theorem.  The hard part is multi-crossing geometry in the
centered annulus:
\begin{itemize}
\item the net signed radial displacement on each ray should be controlled by
the same Steklov spectral gap as in the graph proof;
\item same-ray cancellations, where \(D\) has both an inner missing piece
and an outer excess piece on the same direction, should cost local radial
capacity, plausibly linearly in their cancelled volume;
\item no file in Plan 2 or Plan 3 currently proves this capacitary
decomposition for the torsion energy.
\end{itemize}

Plan 2's slice-rearrangement material supplies the same-center geometric
technology behind \((4.3)\): per-ray good/bad decompositions, bad-ray
multiplicity absorption, oriented radial trace, and estimates of the form
\(|E\Delta B(z)|^2\le C\Pi(E,z)\) under an annular cap.  It does not supply a
single-domain torsion lower bound converting \(E(D)-E(B)\) into one of those
same-center annular defects.  That is the precise remaining bridge.

\section{Soundness verdict}

\begin{itemize}
\item \textbf{Overclaim found.}  Chain B's unweighted Hessian inequality
\((T)\) is false.  The false claim is in the transition
\[
   \delta_T(D)\quad\Longrightarrow\quad
   \int_D|D^2u_D+\tfrac1n I|^2,
\]
not in Step 1, Step 3, or Step 5.
\item \textbf{Underclaim corrected.}  Lack of \(C^1/H^1\)/curvature control
is not by itself fatal in dimension \(2\).  The new radial-graph theorem
shows that small annular amplitude can be enough if one can first prove
single-crossing, or replace single-crossing by an annular-set theorem.
\item \textbf{Centering caveat.}  Any annular theorem must remove the
translation mode.  The fixed-center bound by \(|D\Delta B_1|^2\) is false
without barycenter/optimal-center normalization, as translated disks show.
\item \textbf{Best non-Plan-1 targets.}  Prove the planar annular entry
(AE), then prove the planar annular bridge (AB), or the stronger
radial-moment coercivity (4.4).  This would use the correct
\(H^{1/2}\)/Fraenkel scale and would not rely on BDV.
\item \textbf{What not to pursue.}  Do not try to close Plan 3 through a
single-level Hessian square function, weighted or unweighted.  Its spectral
order is too high.
\end{itemize}

\end{document}
